\section{Probability references and the reach of the method}
\label{p1:discussion}

The main theorem gives a probability reference with the exact
first two factorial cumulants. A second positive reference has
an elementary sampling interpretation. It belongs to the classical
Hermite family, whose single and double jump representation is
studied by Kemp and Kemp~\cite{kempkemp1965}.

The same stopped expansion also retains the unsigned integer
positions. Lemma~\ref{p1:limits:vector-reference} compares
\((J_x)_{x\in I_N}\) with a compound-Poisson field built from
independent single arrivals at sites and double arrivals along
good arithmetic pairs. Its remainder is
\(O(N^{-4/5+\varepsilon})\), including every fixed
polynomial weight in the total count.
Section~\ref{p1:limits:positions} develops that comparison and
uses it to justify pair deletion. The target keeps the locations
of both endpoints, but does not record run signs or exact excess
lengths.

\begin{corollary}[Hermite approximation with single and double jumps]
\label{p1:cor:compound}
For all sufficiently large \(N\), let \(U_N,V_N\) be independent
Poisson variables with means \(\mu_N-\alpha_N\) and
\(\alpha_N/2\), respectively. Then
\[
 \TV\bigl(\Law(Z_N),\Law(U_N+2V_N)\bigr)
       \ll_{C,\varepsilon}N^{-4/3+\varepsilon}.
\]
For every fixed \(h\ge3\), the entirely explicit reference
\[
 \widehat U_N+2\widehat V_N,\qquad
 \widehat U_N\sim\Pois(\lambda_N-\widehat\alpha_N^{(\le h)}),\qquad
 \widehat V_N\sim\Pois(\widehat\alpha_N^{(\le h)}/2),
\]
with independent summands has optimal power \(h/(h+1)\)
and the same leading distance equivalent as the corresponding
quadratic reference in Theorem~\ref{p13:channels:hierarchy}.
In particular the two-channel choice \(h=4\) has
error \(O(N^{-4/5+\varepsilon})\).
No matching lower bound is asserted here for the
exact-parameter Hermite reference.
\end{corollary}

\begin{proof}
Both means are nonnegative for large \(N\).
In the convolution algebra, the law of \(U_N+2V_N\) is
\[
 \exp_*\left(\mu_N\Delta+\frac{\alpha_N}{2}\Delta^{*2}\right).
\]
Expanding the second exponential to first order gives
\(Q_{\mu_N,\alpha_N}\) with an \(\ell^1\) error
\(O_C(\alpha_N^2)\). Indeed \(\|\Delta^{*2}\|_1=4\),
and the exponential remainder is bounded by
\(\frac12(2\alpha_N)^2e^{2\alpha_N}\).
Since \(\alpha_N^2=N^{-4/3+o(1)}\), the main theorem
proves the claim.
For each fixed-height explicit reference, the coefficient
is positive and tends to zero uniformly in the critical band, while
\(\lambda_N\) stays bounded away from zero. Thus both proposed
Poisson means are nonnegative for all sufficiently large \(N\).
Repeat the same expansion with parameters
\((\lambda_N,\widehat\alpha_N^{(\le h)})\), and apply
Theorem~\ref{p13:channels:hierarchy}. The additional norm
error is \(O_C((\widehat\alpha_N^{(\le h)})^2)
=N^{-4/3+o(1)}\), smaller than the positive first omitted
channel term of order \(LN^{-h/(h+1)}\).
The reverse triangle inequality preserves its equivalent
and hence its optimal exponent.
\end{proof}

The double jump represents the leading pair interaction at
the level of the scalar law. The contribution here is the
precision of the Hermite approximation for \(Z_N\).
Its double jump intensity uses either the exact second
factorial cumulant or a complete finite channel sum.
This is a representation of the reference law; it does not
decompose the actual run starts into independent components.
Keeping only the leading equivalent of that coefficient
does not preserve the precision of a fixed-height truncation.
The exact Hermite law already contains the quartic term
\(\alpha_N^2\Pi_{\mu_N}*\Delta^{*4}/8\), so the lower
bound for the quadratic reference does not transfer to it.

This distinction is already present in the locally dependent
setting of \v{C}ekanavi\v{c}ius and Vellaisamy
\cite[Theorem~3.1, (3.5), and Theorem~3.2, (3.6)]{cekanaviciusvellaisamy2015}.
Their quadratic reference is
\(\Pi_{\Gamma_1}*(\delta_0+\Gamma_2\Delta^{*2})\), whereas
the exponential reference is
\(\exp_*(\Gamma_1\Delta+\Gamma_2\Delta^{*2})\), with
\(\Gamma_2=(\Var Z-\E Z)/2\).
The quadratic estimate contains an additional squared error term
that is absent from the exponential estimate.
These results concern sums reduced to one-dependence under
specified moment conditions. We use this comparison to place the
references in context; the arithmetic estimates here provide
their approximation to the multiplicative run count.

\subsection{Computing the correction coefficient}
\label{p1:discussion:charlier}

The coefficient in the total variation equivalent can be
evaluated from four Poisson masses.  We record the elementary
Poisson--Charlier calculation, together with a bound for every
order. These are standard identities for signed Poisson
approximations. The orthogonality bound and the exact
second-difference norm appear in Zacharovas and
Hwang~\cite[(2.13), Remark~4.4]{zacharovashwang2010}; see also
\cite{bkn2014} for the Poisson--Charlier expansion framework.
We include the calculation to fix the normalization of
\(\Delta\), the endpoint conventions, and the factor relating
the \(\ell^1\) norm to total variation.

\begin{proposition}[Norms of Poisson differences]
\label{p1:discussion:charlier-norm}
For $\mu>0$, put $\pi_\mu(k)=e^{-\mu}\mu^k/k!$ when $k\ge0$
and $\pi_\mu(k)=0$ when $k<0$.  Define
\[
 u_\pm=\mu+\frac12\pm\frac{\sqrt{4\mu+1}}2,\qquad
 a=\lfloor u_-\rfloor+1,\qquad b=\lceil u_+\rceil-1.
\]
Then
\begin{equation}
 \|\Pi_\mu*\Delta^{*2}\|_1
 =2\bigl\{\pi_\mu(a-1)-\pi_\mu(a-2)
             +\pi_\mu(b-1)-\pi_\mu(b)\bigr\}.
 \label{p1:discussion:finite-norm}
\end{equation}
More generally, for every integer $r\ge0$,
\begin{equation}
 \|\Pi_\mu*\Delta^{*r}\|_1
       \le\min\left\{2^r,\frac{\sqrt{r!}}{\mu^{r/2}}\right\}.
 \label{p1:discussion:general-norm}
\end{equation}
In particular,
\begin{equation}
 \|\Pi_2*\Delta^{*2}\|_1=\frac{10}{3e^2},
 \qquad \mathcal C(2)=\frac{5}{6e^2}.
 \label{p1:discussion:lambda-two}
\end{equation}
\end{proposition}

The proof, including the normalization of the Charlier polynomials and
the telescoping at integer roots, is given in
Section~\ref{p14:comp:poisson-differences}.

The function $\mathcal C$ is continuous and strictly positive
on $(0,\infty)$.  Indeed, coupling two Poisson variables by
an independent increment gives
$\|\Pi_\mu-\Pi_\nu\|_1\le2(1-e^{-|\mu-\nu|})$.
Convolving with $\Delta^{*2}$ proves continuity; the mass
at zero gives positivity, as in
\eqref{p1:close:constant-bounds}.  Thus the finite formula
remains continuous when one of its two roots crosses an integer.

For completeness, let
$H_{\mu,\theta}=\Law(\Pois(\mu-\theta)+2\Pois(\theta/2))$,
with independent summands and $0\le\theta<\mu$.
The expansion of its convolution exponential gives the
parameter-uniform estimate
\begin{equation}
 \left|\TV(H_{\mu,\theta},\Pi_\mu)
                     -\theta\mathcal C(\mu)\right|
       \le\theta^2e^{2\theta}.
 \label{p1:discussion:compound-coefficient}
\end{equation}
To verify the constant, write
$H_{\mu,\theta}
 =\Pi_\mu*\exp_*(\theta\Delta^{*2}/2)$.
The $\ell^1$ norm of the exponent is at most $2\theta$,
so the remainder after its linear term has norm at most
$2\theta^2e^{2\theta}$.
The reverse triangle inequality and
$\TV=\|\cdot\|_1/2$ prove
\eqref{p1:discussion:compound-coefficient}.
In particular, the factor $1/4$ in $\mathcal C$ consists
of the factor $1/2$ in the quadratic correction and the
factor $1/2$ in the total variation convention.
Neither \eqref{p1:discussion:general-norm} nor the remainder
constant in \eqref{p1:discussion:compound-coefficient} is
asserted to be optimal.

\subsection{A transferable proof structure}

The method starts by choosing a reference that already contains the
frequent local effects. Here the exact Bernoulli marginals absorb repeated
indices and the retained covariance term absorbs isolated pairs. The
choice is part of the proof: replacing it by a Poisson law too early would
charge the same diagonal effects again.

Next one retains a hereditary family of configurations with explicit joint
probabilities. The specialization of the abstract lace expansion
\cite{zeilberger1997} in Section~\ref{p1:stop:section}
stops at its first excluded prefix.
It leaves the later variables in positive conditional kernels. The
probabilistic task is therefore to bound the total mass of that boundary,
while expanding incompatible configurations with their signs before taking
absolute values.

Three distinct arithmetic controls supply that bound. Conditional
full-support counts give a half power for each new block when the numbers
of sources and targets are fixed. Large rank losses are confined to thin
families, with a pointwise rank bound obtained from short relations and a
common curve. Finally, certificate and saturation estimates sum the
remaining configurations at growing orders. The last step is an independent
uniform estimate; it is not obtained by letting the fixed-order constants
grow. The reconstruction theorem expresses their combined effect as
$N^\varepsilon A_3+N^{-4/3+\varepsilon}$.

In another model, private prime coordinates could be replaced by other
local coordinates that force rarity or yield a sparse dependency graph.
To apply the method one would need counterparts of these conditional
counts, rank bounds, sign cancellations, and uniform tails. The
combinatorial stopping identity is general; the arithmetic estimates
proved here are specific to multiplicative runs.

The lower bound must then test the chosen reference. For the quadratic
truncation the fourth Charlier polynomial detects
$3\alpha_N^2+\kappa_{4,N}$, whereas for the Hermite reference it detects
only $\kappa_{4,N}$. The argument therefore identifies the precision of a
specified approximation, rather than assigning a universal exponent to
dependent rare events.

\subsection{Precision, parameters, and remaining questions}

Table~\ref{p16:tab:references} collects the precisions for the
different references. Beyond those results, determining a full
equivalent for the exact quadratic remainder or the third factorial
cumulant remains open here. So do a sharp lower bound for the exact
Hermite reference and the error after a cubic or quartic correction.
Optimizing over all Poisson means is another question: the equivalents
proved here use the specified means in Corollary~\ref{p1:cor:tv}.
The coefficients determine explicit reference laws, but the constants
in the asymptotic errors and the thresholds in \(N\) are not given
numerically.

The independent proof of bounded fixed moments is also
essential. It precedes the law approximation and permits
the weighted scalar improvement in
Corollary~\ref{p13:optimal:weighted}.
Total variation controls bounded observables,
whereas powers of the count are unbounded. Conversely,
even eventual exact agreement at each fixed moment order, together
with uniformly bounded exponential moments, does not give
a polynomial total variation rate. Section~\ref{p1:moment-obstruction} of the companion
provides a finite-difference example. The distinction between
a fixed order and the logarithmic range used in the stopped
expansion is thus substantive.

The stopped expansion does not require uniformly bounded exponential
moments for the original count. The arithmetic cylinder in
Proposition~\ref{p1:limits:cylinder} explains why such a bound cannot be
imported from a nearby truncated count.

For the positional comparison in
Lemma~\ref{p1:limits:vector-reference}, the individual edge weights
\(q_e\) specify geometry that the scalar coefficient
\(\widehat\alpha_N\) alone cannot recover. A sign-resolved field,
rare doublet count, or finer Palm distribution would require its
own target and error estimate.
